\chapter*{Purpose, scope and reading conventions}
\addcontentsline{toc}{chapter}{Purpose, scope and reading conventions}
These notes organize a single-variable undergraduate analysis course by mathematical
content rather than by copying the chapter order of one textbook. The primary
comparison sources are Zorich I and II, Rudin, Apostol and Pugh. Tao's first
analysis volume supplies a second foundational viewpoint; Olver is used only
for bounded real-variable asymptotic supplements. The source register explains
the selected scope and the original-source locations.

The stable identifiers \texttt{K03.xx.name} refer to the accompanying knowledge
registry. They identify mathematical content, not Lean declarations or a prescribed
number of formalization tasks. A definition, a counterexample, a representation
comparison and a theorem are different kinds of content. The main text retains
the fourteen original core modules and five explicit supplements. Five records
only route material to another direction: Riemann--Lebesgue and Stieltjes--measure
interfaces, advanced asymptotics, general special functions, and general Fourier
transform theory. Those routing records assert no unproved theorem and are not
prerequisites of the proved D03 content.

We use the standard complete ordered real field and its complex extension, along
with elementary set theory and finite arithmetic. These are declared shared
foundations, not newly proposed axioms for a Lean project. Consequences used by
analysis, including completeness arguments and convergence criteria, are proved
in the text. A separate construction of the reals by cuts or rational Cauchy
sequences is not required to use this interface. We do not assume Lebesgue
integration, general multidimensional change of variables, complex contour
integration or a general ODE existence theorem in order to prove one-variable
results.

Unless stated otherwise, an interval is an interval in $\R$, a scalar belongs to
$\R$ or $\C$, and derivatives use real increments. An exact mean-value equality
with one intermediate point is a real-valued statement; complex-valued functions
have a norm estimate instead. The domains of logarithms, real powers and inverse
functions, endpoint continuity, integrability and directions of limits are part
of each assertion. A total integral expression in a future formal library never
replaces an integrability hypothesis.

The elementary functions are built from general series and power-series results;
the logarithm is first the inverse of the exponential. Its integral formula is
then a comparison theorem. For this reason the exposition places general series
before the elementary-function chapter. A few applications refer forward to
later established functions or integration rules; these applications are not
used to prove the foundational series or completeness results. The companion
audit records the prerequisite choices and actual review limits.

The mathematical prose and proofs are newly written syntheses. Source references
identify comparison material; they do not mean that every theorem or exercise in
every book has been included. Distinct useful hypotheses are kept distinct, and
the generality of a comparison bridge is not imposed on the underlying theory.
In particular, refinement and unrestricted mesh versions of a Stieltjes integral
are genuinely different for some jump integrators.
